%!TEX program = lualatex
\documentclass[ngerman,headheight=70pt]{scrartcl}
\usepackage[ngerman]{babel}
\usepackage{blindtext}
\usepackage[utf8]{luainputenc}
\usepackage{scrlayer-scrpage}
\usepackage[T1]{fontenc}
\usepackage{fontspec}
\usepackage{multicol}
\usepackage{enumitem}
\usepackage[german=quotes]{csquotes}
\usepackage[bottom=2cm,footskip=8mm]{geometry}
\setmainfont{Source Sans Pro}
\MakeOuterQuote{"}

\setlength{\parindent}{0pt}
\setlength{\parskip}{6pt}

\pagestyle{plain.scrheadings}
\thispagestyle{scrheadings}
\lohead{\vspace*{-1.5cm}\includegraphics[scale=0.4]{../images/stupa_logo_rot.pdf}}
\rohead{\\Präsidium des\\Studierendenparlaments}

\lofoot[Universität Hamburg · Präsidium des Studierendenparlaments\\\hrule Von-Melle-Park 5 · D-20146 Hamburg]{Universität Hamburg · Präsidium des Studierendenparlaments\\\hrule Von-Melle-Park 5 · D-20146 Hamburg}
\cofoot[]{}
\cohead[\thepage]{}

\renewcommand{\thesubsection}{TOP \arabic{subsection}}
\renewcommand{\thesubsubsection}{\arabic{subsubsection}.}
\setcounter{subsection}{-1}

\begin{document}
    UHH · StuPa-Präsidium · Von-Melle-Park 5 · D-20146 Hamburg

    %%% PAGE ONE (EDIT LATER) %%%%
    % Hier den Titel des Protokolls eintragen
    \section*{Protokoll der 7. Sitzung des Studierendenparlaments vom 14. Juli 2016}

    % Protokollanten eintragen
    \textbf{Protokoll: Jim Martens}\\
    \textbf{Ort: B136 in der ehemaligen HWP}\\
    \textbf{Beginn: 18.22 Uhr}\\
    \textbf{Ende: 19.21 Uhr}

    % Wohl die komplizierteste Stelle im Protokoll
    % Als Faustregel sollten nicht mehr als 3 Namen auf einer Zeile stehen
    % Einfach an dieser Vorlage orientieren
    % Im Idealfall ist hier nicht viel zu ändern, da kaum jemand fehlt.
    % Rücktritte müssen natürlich pro Protokoll angepasst werden
    \vspace{0.5cm}
    \begin{tabular}{ll}
        Anwesend: & \\
            RCDS (5 Sitze): & Ramon Weilinger, Antonia Niecke, Ramin Shakiba, \\
                            & Jennifer Maack \\
             CampusGrün (14 Sitze): & Laura Franzen, Geoffrey Youett, Philipp Droll, \\
                                   & Tahnee Herzig, Svenja Horn, Jim Martens,\\
                                   & Mirzo Ulugbek Khatamov, Armin Günther, Martin Sievert\\
                                   & Elena Rysikova, Yasemin Günther \\
             Bier-Liste (2 Sitze) : & Jakob Pape, Karen Martirosian \\
             WiWi (2 Sitze): & Claas-Friso Hente \\
             Unicorns (5 Sitze): & Johannes Peplow, Annkathrin Löffler, \\
                                 & Andreas Hartkamp, Marielle Hermstrüwer \\
             Liste LINKS (3 Sitze): & Gunhild Berdal, Till Petersen, Sinah Mielich \\
             HWP (2 Sitze): & Ajdina Karahasan \\
             MIN (4 Sitze): & \\
             SDS* (3 Sitze): & Mena Winkler, Jacob Petersein, Artur Brückmann \\
             Bart-LISTE (2 Sitze): & Timo Zeimet, Dominic Laumer \\
             LHG (1 Sitz): & \\
             harte zeiten (1 Sitz): & Tobias Berking \\
             Jura (1 Sitz): & \\
             AL (2 Sitze): & Karima Schulze, Henri Weber \\
            & \\
        Entschuldigt: & Benjamin Welling (RCDS),\\
                      & Melf Johannsen, Mario Moldenhauer (CG), Samet Gunay,\\
                      & Kay Zöllmer (WiWi), Katharina Kucza (Unicorns), Ailina Salten, \\
                      & Lotte Rullkötter, Thea Wahlers (alle MIN), Tobias Heisig (LHG)\\
                      &\\
        Unentschuldigt abwesend: & Jan Detampel (MIN), \\
                                & Johann Baumhoefener (Jura)\\
                                &\\
        Rücktritte: & Freya Schmitz (CG) \rightarrow Sandra Fietz\\
    \end{tabular}
    \newpage
    %%% Vorgeschlagene Tagesordnung %%%
    \underline{Vorgeschlagene Tagesordnung}
    \begin{enumerate}[label={\textbf{Top \theenumi}},leftmargin=*]
        % folgende Zeile bei Sitzungsfortsetzungen einkommentieren und x durch
        % die Zahl des zuletzt behandelten TOPs ersetzen
        %\setcounter{enumi}{x}
        \item Auswertung der Diskussion mit Katharina Fegebank  (45 Min.)
        \item Einstieg ins Wintersemester unter besonderer Berücksichtigung des Hamburger Haushalts (45 Min.)
        \item Zivilklausel/Transparenzklausel (45 Min.)
        \item Etwas Besseres: Feuer und Flamme für soziale, ökologische und friedliche Stadt (V1617-031) (45 Min.)
        \item Unterstützung bundesweiter Friedensdemo (V1617-032) (45 Min.)
        \item Fernmeldeamt für Uni nutzen (V1617-033) (45 Min.)
        \item Verschiedenes (5 Min.)
    \end{enumerate}

    \newpage

    % Hier beginnt Teil A. Bei Fortsetzungen kommt hier nur der Teil über
    % die Feststellung der endgültigen Fassung der Tagesordnung.
    % Für sonstige Sitzungen ist dies bereits das Gerüst für die Formalia.
    % Man muss nur noch die protokollierten Inhalte eintragen.

    {\Large \underline{Teil A}}

    \subsection{Formalia}

    \textit{Ramon eröffnet die Sitzung um 18.22 Uhr.}

    \subsubsection{Geschäftsbericht Präsidium}

    Ramon stellt den vorgeschlagenen Sitzungsplan für das Wintersemester
    vor.

    Die Post wurde kontaktiert, um eine Sonderpostleitzahl zu beantragen.

    \subsubsection{Anfragen an das Präsidium}

    F: Wann ist die Abgabe für Kandidierendenlisten?
    A: 46. KW, 8. Dezember Briefwahlabgabe

    F: Warum nicht schon Sitzung Mitte Oktober?
    A: Damit erste Sitzung nicht schon in der 1. VL-Woche ist.

    \subsubsection{Geschäftsbericht AStA}

    Franziska hält den Bericht.

    Der AStA sei gegen Zwangsexmatrikulation. Dahingehend wird argumentiert, dass
    die Zahlung des Semesterbeitrags Ausdruck des Willens ist weiter studieren zu
    wollen.

    Die Raumverteilung soll wieder in Präsidialverwaltung eingegliedert werden.

    Es soll ein Ausschuss zu Religion im AS gegründet werden. Der AStA hat sich
    mit dem Thema Religion beschäftigt.

    Die AS-Sitzung fand statt. Das Ergebni ist, dass es einen weiteren Dies
    Academicus geben werde. Mena wird etwas zur RPO vortragen.

    Mena erzählt über den aktuellen Stand zur RPO. Der ALSt-Entwurf der RPO
    wird in einzelne Fakultäten und Fachbereiche zur Stellungnahme gegeben.
    FSRe sollen sich beteiligen, damit möglichst weitestgehende RPO
    am Ende steht. Kritische Punkte könnten sein: Anwesenheitspflicht,
    Modulbegrenzung, unbegrenzte Wiederholung von Prüfungen.
    Gegenargument zu unbegrenzter Wiederholung ist häufig eine höhere
    Prüfungslast. Konter: Insgesamt werden Prüfungen zurückgefahren.

    AStA wird erarbeitete Stellungnahmen zusammenstellen. Eine Stellungnahme
    des AStA zum vorliegenden Entwurf soll erarbeitet werden. Möglicherweise
    könnte auch ein Video gemacht werden.

    Im September soll ein größeres gemeinsames Diskussionstreffen stattfinden.

    Droll erzählt über die Veranstaltung zu libertären Mauerkaninchen und
    totalitären Schäferhunden.

    Artur erzählt über einen Artikel im Hamburger Abendblatt, in dem
    Schinnenburg seine große Anfrage zusammenfasst und gegen das Studierendenwerk
    schießt. Einige Leute aus dem Studierendenwerk erklären jetzt, was
    die Aufgabe des Studierendenwerks ist.

    Es gab eine Vertreterversammlung am 6. Juli. 165 Studierende (153 international)
    stehen noch auf Warteliste für Studierendenheime. Einige internationale
    Studierende haben davon erzählt, dass Studierende gewisser Nationalitäten
    nicht mehr aufgenommen werden, weil die Quote für diese Nationalitäten bereits
    voll sei. Der AStA hat zwei Anträge eingebracht. Der Antrag sich im WS
    an Protesten zu beteiligen wurde mit 7:8 abgelehnt. (7 Personen:1 Person Präsidium)

    \subsubsection{Anfragen an den AStA}

    Till fragt, wie es aussieht mit Zugang zu Räumen in den Erziehungswissenschaften
    außerhalb der Öffnungszeiten.

    Droll antwortet: Es gab noch kein Treffen diesbezüglich. AStA hakt noch
    einmal nach.

    Till fragt, wie Auswertung des AStA-Fußballturniers aussieht.
    Eine Refugee-Mannschaft ist auf eigene Kosten von Köln nach Hamburg
    gekommen. Einige Mannschaften kamen direkt von zentralen
    Erstaufnahmen. War eine gute Sache, aber Zelebrierung von Toren
    und Aggressivität war schade. Aber kein Vorwurf an einzelne Personen.
    Sollte jedoch wiederholt werden. Aber es muss ein kulturelles
    Zentrum geben.

    A: Es gab noch keine gesamte Auswertung. Morgen auf der AStA-Sitzung
    wird darüber geredet. Man hätte mehr über die Schiedsrichter sprechen
    müssen und nach welchen Regeln gepfiffen wird. Feedback war, dass es
    einen großen Bedarf gibt und bestenfalls regelmäßig stattfinden sollte.
    Es gab jedoch auch "eine Tonne" Verbesserungsvorschläge für die Veranstalter.

    Till fragt, welche Schlussfolgerungen sich aus der Situation im
    Wissenschaftsbetrieb ziehen lassen. Was hat die Mauerkaninchen-Veranstaltung
    für Erkenntnisse gebracht?

    A: Man kann Spaß haben und kritisch sein. Florian Peters habe viel Zuspruch
    bekommen, auch von Leuten die sonst Anträge schreiben. Es gibt aber keinen
    organisierten zusammenhang von Wissenschaftlern, die solche einzelne
    Kritik aufnehmen und Schlussfolgerungen ziehen könnten. Es gibt viel Kritik,
    aber kaum jemand im wissenschaftlichen System mache etwas. Daher sei viel
    Potential vorhanden Leute wachzurütteln.

    Katharina fragt, wo das Turnier beworben wurde? Vorschlag beim nächsten Mal
    auch Workshops anzubieten.

    A: Der AStA hat auf dem Campus geworben. RefugeeWelcome Karoviertel hat in
    der Stadt geworben. Vielleicht hätte die Resonanz größer sein können. Zunächst
    waren 16 Mannschaften geplant, es sind dann 20 geworden.

    Gunhild fand das Turnier richtig und wichtig und ein gutes Gegenprogramm
    zur EM. Ihrer Meinung hätte es deutlich politischer gestaltet werden können.
    Vorstellung der Mannschaften und gegenseitiges Kennenlernen sollte in
    Zukunft besser gemacht werden.

    Gunhild fragt, ob der AStA dem fzs schon beigetreten ist und ob es erste
    Erkenntnisse diesbezüglich gibt.

    A: Offiziell wird die VS erst im Oktober Mitglied werden. Am ersten WE im August
    ist die nächste Mitgliederversammlung des fzs in Hannover. Der AStA wird teilnehmen.

    Es wird quasi in eine Aktuelle Stunde übergegangen.

    Katharina fände es toll, wenn einmal im Semester ein Raum geschaffen wird,
    um sich mit Geflüchteten auseinanderzusetzen.

    Till ist sich sicher, dass alle Flüchtlingsunterkünfte erreicht wurden.
    5 Unterkünfte waren mit einer eigenen Mannschaft anwesend. AStA kann kein
    Turnier mit mehr Mannschaften wuppen. Till ist sehr für Workshops und dafür,
    dass man solche Turniere mehr auf den Breitensportaspekt des Fußballs konzentriert.
    Es sollte kein Gegenprogramm zum Turnier geben.
    Till findet es bekloppt, wenn Leute, die sich als links bezeichnen,
    Gegner als Lusche bezeichnen. Man könnte eine kulturelle Setzung machen.

    Katharina begründet noch einmal ihre Workshopidee.

    Artur begründet, warum Leute womöglich kein Interesse an Fußball haben.
    Es sollte eine alternative Fußballkultur gelebt werden.

    Tobias findet es wichtig, dass sowohl das Fußballspielen als auch Workshops
    bzw. weitere Aktivitäten in einem kulturellen Zentrum nicht gegenläufig
    sein sollten. Es sollte nicht auf ein Endspiel hinauslaufen. Das Spielen
    sollte weniger athletisch, sondern mehr als Zusammenspielen gelten.

    Katharina fragt, ob jeder eine Medaille bekommen hat? Jeder Spieler,
    der am Ende noch anwesend war, hat eine Medaille bekommen. Die Medaille
    zeigt zwei Männer, wo der eine dem anderen irgendwo hinfasst. Es wird
    nicht sonderlich einfach bei den aktuellen Gegebenheiten in den
    ZEAs ein Breitensportturnier zu organisieren. Till wundert sich, warum
    es von den Sportministern noch keine Liga für Refugee-Mannschaften gibt.
    Der AStA könnte das mal anregen.

    \subsubsection{Dringlichkeitsanträge des AStA}

    Keine.

    \subsubsection{Aktuelle Stunde (falls entsprechender Antrag vorliegt)}

    Keine.

    \subsubsection{Feststellung der endgültigen Fassung des Teils B der Tagesordnung}

    \textit{Mit 19 ParlamentarierInnen nicht beschlussfähig. Sitzung wird damit
    offiziell abgebrochen.}

    \textit{Sitzung um 19.21 Uhr beendet.}

\end{document}
